% Figure: how a logarithmic axis is gridded, and how the same data
% point reads differently against a linear and a log axis. The left
% panel shows the uneven minor gridlines of a single decade; the
% right panel contrasts a linear and a log vertical axis carrying the
% same three data values.
\begin{figure}[htbp]
\centering
\resizebox{\linewidth}{!}{%
\begin{tikzpicture}[scale=1.0,
  ax/.style={->, draw=engblack, thick},
  major/.style={draw=engblack, thick},
  minor/.style={draw=engblack!45, thin},
  pt/.style={circle, fill=engred, inner sep=1.4pt},
  ann/.style={font=\footnotesize, align=center},
]
% --- Left: anatomy of one decade on a log axis ---
\begin{scope}[xshift=0cm]
  \draw[ax] (0, 0) -- (0, 6.2) node[left, ann] {log axis};
  % decade from 10^0 to 10^1: minor lines at log10(2..9)
  \foreach \v in {1,2,3,4,5,6,7,8,9} {
    \pgfmathsetmacro{\yy}{4*ln(\v)/ln(10)}
    \draw[minor] (-0.08, \yy) -- (0.4, \yy);
    \node[ann, right, xshift=2pt] at (0.4, \yy) {\v};
  }
  % major at 10^0, 10^1, and partial 10^2
  \draw[major] (-0.15, 0) -- (0.5, 0);
  \node[ann, left, xshift=-2pt] at (-0.15, 0) {$10^{0}$};
  \draw[major] (-0.15, 4) -- (0.5, 4);
  \node[ann, left, xshift=-2pt] at (-0.15, 4) {$10^{1}$};
  \foreach \v in {2,3,4,5,6,7,8,9} {
    \pgfmathsetmacro{\yy}{4 + 4*ln(\v)/ln(10)}
    \ifdim \yy pt < 6.0pt
      \draw[minor] (-0.08, \yy) -- (0.4, \yy);
    \fi
  }
  \node[ann, anchor=south] at (0.6, 6.0) {(a) one decade,};
  \node[ann, anchor=south] at (0.6, 5.6) {minor gridlines};
\end{scope}

% --- Right: same data on linear vs log vertical axis ---
\begin{scope}[xshift=6cm]
  % linear axis
  \draw[ax] (0, 0) -- (0, 6.2) node[left, ann] {linear};
  \foreach \v/\yy in {0/0, 200/1.2, 400/2.4, 600/3.6, 800/4.8, 1000/6.0} {
    \draw[minor] (-0.08, \yy) -- (0.08, \yy);
    \node[ann, left, xshift=-2pt] at (-0.08, \yy) {\v};
  }
  % three data points at 5, 50, 950 on linear scale (scaled /1000*6)
  \foreach \v in {5, 50, 950} {
    \pgfmathsetmacro{\yy}{\v/1000*6}
    \node[pt] at (0, \yy) {};
  }
  \node[ann, anchor=south] at (0.3, 6.0) {(b) linear};

  % log axis
  \draw[ax] (4, 0) -- (4, 6.2) node[left, ann] {log};
  \foreach \e/\yy in {0/0, 1/2, 2/4, 3/6} {
    \draw[minor] (3.92, \yy) -- (4.08, \yy);
    \node[ann, left, xshift=-2pt] at (3.92, \yy) {$10^{\e}$};
  }
  % same three data points 5, 50, 950 on log scale (log10(v)*2)
  \foreach \v in {5, 50, 950} {
    \pgfmathsetmacro{\yy}{2*ln(\v)/ln(10)}
    \node[pt] at (4, \yy) {};
  }
  \node[ann, anchor=south] at (4.3, 6.0) {(c) log};
\end{scope}
\end{tikzpicture}%
}%
\caption[Anatomy of a log axis and the same data on two scales]{Panel
(a) is the anatomy of a logarithmic axis: within each decade the
minor gridlines sit at $2, 3, \ldots, 9$, crowded toward the top of
the decade because the gaps shrink as $\log_{10}$ of consecutive
integers. The signature of a log axis is this uneven crowding, not
the even spacing of a linear axis. Panels (b) and (c) carry the same
three data values, $5$, $50$, and $950$, on a linear axis and a log
axis. On the linear axis the two small values pile up
indistinguishably near the origin and only the largest is readable;
on the log axis all three separate cleanly. The choice of axis
decides which structure the eye can see.}
\label{fig:vol02:ch02:log-axis-reading}
\end{figure}
